\chapter{DOGMA state pools and memory ownership}
\label{ch:evo32}
\section{The address is not the state}
Chapter 31 defined a tenant/request-owned state bound to artifact content,
position and a hidden vector. Reusing physical storage introduces a different
problem. Request A writes slot zero, finishes, and request B receives that
slot. A stale reference from A still names the same address. Address equality
does not mean logical-state equality.

A safe design needs both temporal identity and a reclamation rule. Generation
numbers reject old handles after reuse. They do not, by themselves, prevent an
already-running kernel from touching memory before it has finished. This chapter
therefore adds explicit in-flight pins and deferred release to the incoming
package's generational pool.

The original implementation is a serialized CPU inference oracle using copied
scalar snapshots. It is not a concurrent GPU allocator, full DOGMA state layout,
authentication mechanism or secure-erasure implementation. We use it to state
and test the obligations that those systems would have to refine.

\section{Map logical identity to bounded storage}
A pool has capacity $C$, hidden width $D$ and one artifact digest. Every occupied
slot retains owner $(\text{tenant},\text{request})$, consumed-token position,
hidden payload, slot generation, update version and in-flight pin records.
The pool instance has a unique identity separate from its display name.

A \Term{generational handle} contains
\[
 \eta=(\text{pool instance},i,g,\text{owner},\text{artifact}),
 \qquad 0\le i<C.
\]
A lookup validates the caller's owner pair, handle fields and current slot
metadata before exposing a snapshot. Reading with a handle carrying owner A
while supplying caller B is rejected: owner metadata in a handle is not itself
caller authorization.
\Plate{01-logical-physical}{Logical state and physical residency are separate identities. The slot can be reused only under a new generation after reclamation.}{fig:evo32-logical}
\Listing{Pool._slot}

Exact identity checks within trusted Python code are not authentication.
A malicious caller controlling the process could inspect private fields or
forge objects. Production authorization must derive caller identity from a
trusted boundary, not accept it from arbitrary request input.

\section{Type the pool before allocating}
The pool is bound to the artifact's digest and declared width, with an
inference payload of finite float64-compatible scalars and nonnegative integer
position. Boolean dimensions, Boolean positions, nonfinite values and wrong
widths are rejected. Validating a complete payload before publication prevents
a failed initialization from consuming a slot.

One pool class does not imply that full DOGMA has one vector. A real state may
contain multiple banks, finite windows or other carried fields. Their layouts,
dtypes and semantics must be part of the typed artifact contract. Matching
the shape of one tensor does not make two architectures interchangeable.

The reference uses monotonically growing Python integers for generations.
A bounded hardware counter would require an explicit wraparound policy.
Allowing an ancient handle to regain validity after counter wrap would undo
the identity guarantee.

\section{Allocation has an all-or-nothing publication point}
Allocation first validates owner and optional snapshot, then selects a free
slot and publishes its payload/owner/header. Exhaustion raises a capacity
error instead of borrowing another live request's state.
\Listing{Pool.allocate}

The single-threaded reference can treat owner installation as its publication
point. A concurrent implementation needs synchronization covering reservation,
initialization and handle publication. If initialization fails after reservation,
the implementation must scrub and return the slot without publishing an
incomplete handle. Our oracle validates before changing slot metadata.

A pool instance identifier also prevents two different pools with the same
display name from accepting each other's handles. This is a correctness
namespace, not a secret credential.

\section{Generation and version solve different stale-state problems}
Generation changes when a slot is reclaimed. Update version changes when a
live slot's payload changes. They answer different questions:
\begin{enumerate}
\item Does the handle still refer to this allocation? Check generation.
\item Was the result computed from the current live payload? Check version.
\end{enumerate}

Two computations can pin version zero of the same allocation. If the first
installs version one, the second must not overwrite it with a result computed
from version zero. Both handles have the same valid generation; generation
alone cannot reject the lost update.
\Plate{03-handle}{Two fences: allocation generation rejects address reuse; live update version rejects a result based on obsolete state.}{fig:evo32-handle}

Tagged pointer/counter structures occur in established concurrent algorithms
\cite{evo32-michael}. Our pool is not an implementation of their queue, nor a
claim of lock-free progress. Tagged identity and safe memory reclamation remain
separate design obligations.

\section{Pin in-flight work before release}
A \Term{state pin} captures a handle, unique completion serial, live update
version and immutable input snapshot. The slot records that pin until its
completion or explicit discard is acknowledged. Pinning a retiring allocation
is forbidden.

Calling release on an unpinned allocation scrubs its API-visible payload,
clears its owner and advances the generation. Calling release while pinned
marks the slot retiring but does not expose it as free. Ordinary read, write
and new-pin operations then fail. Existing pin completions may settle, but
their candidate results cannot be installed.

This lifecycle prevents a cancellation request from being mistaken for a
kernel-completion signal.
\Plate{02-lifecycle}{Deferred reclamation. A retiring slot remains occupied until the last in-flight acknowledgement; only then may it be scrubbed and reused.}{fig:evo32-lifecycle}

\section{Fence completion and settle every pin}
The completion path validates the exact pin object, its allocation, caller
and payload. A candidate may install only when the slot is not retiring and
its live version matches the captured input version. Its position must be the
pinned position plus one live token.
\Listing{Pool.complete}

A discarded completion still removes its pin. When a retiring slot loses its
last pin, reclamation becomes legal. Duplicate acknowledgement is an error.
An invalid candidate leaves the pin live so that a caller can explicitly discard
it through the same completion API; silently dropping the pin would hide
unfinished-work accounting.

The pool does not verify that arbitrary candidate values actually satisfy the
model recurrence. The trusted executor must produce those values using
Chapter 31's native semantics. The pool verifies ownership, layout, temporal
identity and update boundary, not mathematical reachability of a hidden vector.

\subsection{A cancellation trace with explicit capacity}
For one slot initialized at position three:
\begin{center}\input{\Generated/timeline.tex}\end{center}
Release after pinning leaves zero free slots. The late result is rejected,
but acknowledging it removes the last pin and permits scrubbing.
The next owner receives zero state under a new generation. The old handle
cannot access the new allocation even though the index is unchanged.

This trace is a deterministic event ordering, not a concurrent stress result
or GPU latency measurement. A leaked pin causes retained capacity; it must be
diagnosed rather than bypassed by forcibly reassigning the live slot.

\section{Clone values without sharing writable storage}
The conservative clone reads an immutable snapshot and allocates a new slot
for a new owner in the same tenant. Later writes to the clone do not affect
the source. Cross-tenant clone is rejected because no export authorization
policy is implemented.
\Plate{04-clone}{Inference clone preserves values and token boundary but allocates independent logical ownership. Same-tenant copying is not permission to export another tenant's state.}{fig:evo32-clone}
\Listing{Pool.clone}

The reference copies scalars on input and output. Mutating a caller-provided
list after allocation cannot mutate the slot. A returned snapshot may survive
release, because it is an independent value, not a writable view of reused
storage. Clone-by-copy is the semantic oracle for later copy-on-write work.

Clone costs an additional occupied slot. On exhaustion, it fails without
changing the source. Identical prefixes alone do not justify mutable aliasing.

\section{Bridge to the actual Chapter 31 native state}
The implementation imports Chapter 31's reference rather than inventing a
second recurrence. A fresh native tensor is constructed from the pool snapshot,
with identical owner, artifact, position and hidden values.
\Listing{to_native}

Saving a native state validates it against its model and owner, checks the
pool's artifact digest, then copies its values. A separately materialized
native tensor cannot change the pool through in-place mutation. Tests consume
a real fixture token, save the resulting state and verify exact values on
reload.

This is an \Term{inference snapshot}, not a differentiable training continuation.
Converting tensors to scalar data and rebuilding them removes the graph.
Chapter 31's derivative parity is not preserved by this pool serialization,
and we do not claim otherwise. A training-state allocator needs a distinct
graph-lifetime contract.

The bridge's ordinary save path requires no live pins; asynchronous candidates
must instead use the fenced completion API. Combining the pool with the shell's
output publication requires a joint serialized policy for both resources.
This chapter does not claim that two independent APIs compose into a durable
transaction automatically.

\section{Workspace has a shorter lifetime}
A \Term{step workspace} holds temporary projections or reductions, while the
persistent state must survive continuation boundaries. The teaching workspace
returns immutable copied values, validates shape and finiteness, and refuses
use after close. Closing is idempotent and resets its API-visible payload.
\Plate{05-workspace}{Persistent state and step-local scratch have different lifetimes. Captured input/output snapshots are copied values, not borrowed writable addresses.}{fig:evo32-workspace}

An actual GPU workspace may be borrowed by an asynchronous kernel. Its reuse
must wait for the corresponding completion event or synchronization, just as
persistent state reclamation waits for pins. A lexical function return is not
necessarily a device-completion boundary.

\section{Scrubbing: specify the threat model honestly}
The reference replaces payload with zeros and clears owner/position metadata
before exposing a free slot. This prevents its ordinary API from returning a
previous request's values after reuse.

It does not securely overwrite every copy in Python memory. Old immutable
snapshots can remain referenced, and garbage collection is not cryptographic
erasure. Nor does it overwrite VRAM, clear driver caches or protect crash dumps.
A production erasure claim requires its own platform-specific threat model and
verification. Resetting logical position alone is still insufficient: every
persistent carried field must be initialized consistently before reuse.

\section{Account for capacity and physical costs separately}
At every event,
\[
 C=C_{\rm free}+C_{\rm owned}+C_{\rm retiring}.
\]
Retiring slots count as occupied even if their request is terminal.
Pins are counted separately because multiple pins may refer to one occupied
slot.
\Plate{06-isolation}{Canary reuse and conservation of capacity. Terminal requests may still hold retiring slots until work settles.}{fig:evo32-isolation}

For a contiguous float64 arena, payload size would be $8CD$ bytes.
The reference's \texttt{payload\_bytes} reports this logical scalar-payload
quantity, not Python heap usage. Actual memory also includes metadata, object
overheads, retained snapshots, kernel buffers, allocator reservations and copies.
If cloning every branch, those extra occupied slots must enter admission policy.

A useful metric panel reports free/owned/retiring counts, outstanding pins,
failed allocations, stale-handle events and reclamation delay. It should not
dump hidden values or prompt-derived state into routine telemetry.

\section{Test against an independent ownership ledger}
The tests maintain a separate map of owners to expected snapshots while
performing deterministic randomized allocation/release sequences. Distinct
canaries identify each request. After reuse, recently stale handles must fail
and all new zero-initialized allocations must satisfy their initialization
contract.

Targeted cases add deferred cancellation, multiple outstanding pins, duplicate
completion, live version conflict, malformed allocation, clone exhaustion,
caller mismatch, artifact mismatch and native-bridge storage independence.
These tests are serialized: no race, device stream ordering, synchronization
primitive or production performance is claimed.

For a concurrent implementation, show that each operation can be linearized
to a permitted sequential outcome. Force allocate/release, pin/cancel,
clone/release and completion/reuse races, rather than assuming a high-volume
throughput test covers those interleavings.

\section{Twenty worked exercises}
\begin{enumerate}
\item \textbf{Identity.} Is slot zero a stable request identity?
\emph{Answer:} no; reuse changes its logical owner.
\item \textbf{Generation.} What rejects a pre-release handle?
\emph{Answer:} its old generation differs from current metadata.
\item \textbf{Version.} What rejects a second result from obsolete live state?
\emph{Answer:} its pinned update version differs from the current version.
\item \textbf{Namespace.} Can equal pool display names alias?
\emph{Answer:} not here; pool-instance identity is distinct.
\item \textbf{Caller.} Does a handle's owner field authenticate the caller?
\emph{Answer:} no; caller identity requires a trusted authorization boundary.
\item \textbf{Allocation failure.} An input snapshot has NaN.
\emph{Answer:} validation fails before any slot is published.
\item \textbf{Capacity.} The only slot is retiring with one pin.
\emph{Answer:} there are zero free slots; allocation must fail or queue externally.
\item \textbf{Cancellation.} Can release scrub a pinned slot immediately?
\emph{Answer:} no; mark retiring and wait for acknowledgements.
\item \textbf{Late result.} May a retiring slot install a candidate?
\emph{Answer:} no; completion settles its pin but discards the candidate.
\item \textbf{Two pins.} The first settles after release.
\emph{Answer:} retain the slot until the second also settles.
\item \textbf{Duplicate acknowledgement.} Can the same pin settle twice?
\emph{Answer:} no; it is absent after the first completion.
\item \textbf{Lost update.} Two pins capture version zero; one installs.
\emph{Answer:} the second result is discarded because version is now one.
\item \textbf{Position.} A pin captured position three. Is candidate five legal?
\emph{Answer:} not for this one-live-token completion contract.
\item \textbf{Clone.} Does modifying a branch alter its source?
\emph{Answer:} no; storage is independent.
\item \textbf{Export.} Does clone allow a different tenant?
\emph{Answer:} no; an explicit export policy is absent.
\item \textbf{Snapshot lifetime.} Can a copied snapshot outlive release?
\emph{Answer:} yes; it is an independent value, not a handle to reused storage.
\item \textbf{Training.} Does scalar snapshot reload preserve autograd?
\emph{Answer:} no; this is an inference bridge.
\item \textbf{Memory.} For $C=4,D=2$, what is logical float64 payload size?
\emph{Answer:} $64$ bytes, excluding metadata and Python overhead.
\item \textbf{Erasure.} Does replacement with zeros prove secure erase?
\emph{Answer:} no; copies and platform storage require separate verification.
\item \textbf{Concurrency.} Do these serialized tests prove GPU race freedom?
\emph{Answer:} no; a device implementation must refine the contract and test
real completion/reclamation races.
\end{enumerate}

\section{Handoff to checkpoints and shared prefixes}
Chapter 33's checkpoint, branching and prefix-state candidate must restore
fresh ownership rather than revive a stale pool handle. Its optimized sharing
must refine independent-copy semantics and retain model, semantic version and
exact token boundary. Integrity, authorization and graph lifetime remain
separate obligations. The bounded pool here supplies the ownership and
reclamation foundation for that next chapter.
